\section{Conclusions}
\label{sec:conclusions}

Read every day, the fixed-window tests of VaR back-testing reject correct models far more often than their nominal size: the Kupiec test on the last 250 days in \nNullSizeKupTrail\,\% of 1,000-day series and the traffic light in \nNullSizeLight\,\%. Mixtures of likelihood-ratio e-processes over a grid of alternatives fixed before the data keep the size at any stopping time and measured at \nNullSizeE\,\% for the VaR hit rate (\nNullSizeELo\ to \nNullSizeEHi\,\%), \nNullSizeCluster\,\% for hit clustering and \nNullSizeBattery\,\% for their average, with a proof that rests on the martingale property and Ville's inequality.
The monitors detect a hit rate four times too high at a median of \nSeedsFourMed\ days (the upper-side monitor \nSeedsUpFourMed), against the 250 days the traffic light needs, and a hit rate three times too high at \nSeedsThreeMed\ days, which is not below the window. The upper-side monitor, a variant chosen after the two-sided results, has \nSeedsUpFourMed\ and \nSeedsUpThreeMed\ days and needs a prospective confirmation. The price in power at day 250 is \nSeedsThreeCost\ points at three times the nominal level and \nSeedsFourCost\ at four times for the two-sided monitor, and \nSeedsUpThreeCost\ and \nSeedsUpFourCost\ for the upper side;
the first is above the 10 points allowed by the hypothesis and the three others are under it. The delay and the power-cost hypotheses therefore fail at three times the nominal rate and hold at four. Any monitor valid at every time has a power at day 250 bounded by the most powerful test of the size it has spent, \nExactThreeNpE\,\% at three times the nominal level.

On real series, the e-process flags all 18 models in the equity crisis at a median of \nCElevenEquityNinetyEMed\ days, flags \nCElevenCalmNinetyEFlag\ models in the calm control window where the traffic light read daily flags \nCElevenCalmNinetyLightFlag, and rejects two cohorts of small-business loans at 12 months where the five-year light is red 48 months later. For the defaults of rating grades the monitor keeps its size when defaults are binomial and loses it when the defaults of a year share a systematic factor, rejecting a correct PD in \nRhoZeroThreeAvg\,\% of the histories at an asset correlation of $0.03$, and it detects a 5\,\% relative error in almost no history where a test read once detects \nPdOneZeroFiveChiOne\,\%. At the size the monitor measured (\nPdcMatchedSize\,\%) a chi-square rule read every year detects that error in \nPdcOneZeroFiveChiMatched\,\% of the histories (\nPdcOneZeroFiveChiSim\,\% at a nominal size of 5\,\%), so the defaults are not detected earlier by the monitor. The Vasicek-corrected binomial test of the engine, given the true correlation, keeps its size at every correlation tested (\nPdcZeroThreeVasTrueEvery\ to \nPdcZeroVasTrueEvery\,\%), but with a correlation of \nPdcAssumedRho\ assumed it rejects none of the histories of a PD that is $5\,\%$ to $25\,\%$ too low (\nPdcOneTwoFiveVasAssumedEvery\,\% at $1.25$), so the remedy holds only where the correlation is known.

The next steps are three. A prospective run with fresh seeds should confirm the upper-side configuration. A monitor of the defaults of a grade on a null with a stated asset correlation, built on a beta-binomial or a Vasicek mixture of the likelihood, would remove the limit that the experiments expose. And the same construction can be applied to the e-statistics of expected shortfall of \citet{wang2022}, with the grid of alternatives and the combination rules of this paper.
